% =========================================================================
% PRACTICE EXAM 1: DETAILED SOLUTIONS & SCORING RUBRICS
% =========================================================================
\chapter{Practice Exam 1: Detailed Solutions \& Grading Rubrics}

\section{Section I: Multiple-Choice Detailed Explanations}

\begin{enumerate}[leftmargin=1.5em]
  \item \textbf{Answer: (B)} \\
  \textbf{Explanation:} The $z$-score formula is $z = \frac{x - \mu}{\sigma}$. Substituting the given values:
  \[ z = \frac{91 - 75}{8} = \frac{16}{8} = 2.00 \]
  A $z$-score of 2.00 signifies that an exam score of 91 lies exactly 2.00 standard deviations above the class mean.
  \textit{Distractor analysis:} (A) $1.50$ results from dividing 16 by 10 or calculating $(91-75)/10.67$. (C) $2.25$ is $(93-75)/8$. (D) $1.75$ is $(89-75)/8$. (E) $1.25$ is $(85-75)/8$.

  \item \textbf{Answer: (B)} \\
  \textbf{Explanation:} A categorical (qualitative) variable places an individual into one of several groups or categories. Blood type (A, B, AB, O) is defined by qualitative categories with no meaningful numerical arithmetic. Annual income, blood pressure, commute time, and GPA are all quantitative variables measured on continuous or numerical scales.

  \item \textbf{Answer: (B)} \\
  \textbf{Explanation:} The coefficient of determination is $r^2$. Squaring the correlation coefficient gives:
  \[ r^2 = (-0.92)^2 = 0.8464 \approx 0.846 \]
  By mathematical definition, $r^2$ is always non-negative ($0 \le r^2 \le 1$), so option (A) $-0.846$ is impossible. Interpretation: approximately 84.6\% of the variation in the response variable is explained by the linear relationship with the explanatory variable.

  \item \textbf{Answer: (C)} \\
  \textbf{Explanation:} In a modified boxplot, individual outliers are plotted as separate dots, asterisks, or points beyond the fences. Whiskers do not extend indefinitely to the fences themselves; rather, they extend to the most extreme actual observations in the sample that fall inside the lower fence $Q_1 - 1.5(\text{IQR})$ and upper fence $Q_3 + 1.5(\text{IQR})$.

  \item \textbf{Answer: (B)} \\
  \textbf{Explanation:} Observational studies observe individuals and record variables of interest without imposing treatments. Because no random assignment was utilized, lurking and confounding variables (e.g., diet, exercise, socioeconomic status, sleep habits) cannot be controlled. Thus, observational studies cannot establish cause-and-effect relationships.

  \item \textbf{Answer: (B)} \\
  \textbf{Explanation:} By the general addition rule: $P(A \cup B) = P(A) + P(B) - P(A \cap B)$. Because events $A$ and $B$ are independent, $P(A \cap B) = P(A) \times P(B) = (0.5)(0.4) = 0.20$.
  Therefore:
  \[ P(A \cup B) = 0.50 + 0.40 - 0.20 = 0.70 \]
  \textit{Distractor analysis:} Option (A) 0.90 incorrectly assumes the events are mutually exclusive (disjoint) and fails to subtract $P(A \cap B)$.

  \item \textbf{Answer: (A)} \\
  \textbf{Explanation:} For a binomial distribution $X \sim B(n, p)$:
  \[ \mu = np = 50(0.20) = 10 \]
  \[ \sigma = \sqrt{np(1-p)} = \sqrt{50(0.20)(0.80)} = \sqrt{8.00} \approx 2.828 \approx 2.83 \]
  Option (A) is correct.

  \item \textbf{Answer: (C)} \\
  \textbf{Explanation:} The Central Limit Theorem (CLT) states that for any population distribution with finite mean $\mu$ and finite standard deviation $\sigma$, the sampling distribution of the sample mean $\bar{x}$ becomes approximately normal as the sample size $n$ becomes sufficiently large ($n \ge 30$). The CLT does not alter the shape of the underlying population or individual samples.

  \item \textbf{Answer: (B)} \\
  \textbf{Explanation:} For a right-tailed alternative hypothesis $H_a: p > 0.50$ with test statistic $z = 2.33$, the $p$-value is the area under the standard normal curve to the right of 2.33:
  \[ p\text{-value} = P(Z \ge 2.33) = 1 - P(Z < 2.33) = 1 - 0.9901 = 0.0099 \]
  Using TI-84: \texttt{normalcdf(2.33, 9999, 0, 1)} $= 0.009903$.

  \item \textbf{Answer: (B)} \\
  \textbf{Explanation:} With $n = 16$, the degrees of freedom are $df = n - 1 = 15$. For a 95\% confidence interval, the area in each tail is $(1 - 0.95)/2 = 0.025$. Looking up $df = 15$ with central area 0.95 gives $t^* = 2.131$ (or on TI-84: \texttt{invT(0.975, 15)} $= 2.1314$). Critical value 1.960 corresponds to the standard normal $z^*$ distribution.

  \item \textbf{Answer: (C)} \\
  \textbf{Explanation:} A Chi-Square Test for Homogeneity evaluates whether the distribution of a single categorical variable is identical across two or more independent populations or treatment groups (here, 3 hospital networks). A test of independence uses a single sample categorized by two variables. Goodness-of-fit evaluates a single sample against a hypothesized distribution.

  \item \textbf{Answer: (C)} \\
  \textbf{Explanation:} In simple linear regression inference for the slope $\beta$, degrees of freedom are given by $df = n - 2$ because two parameters (the intercept $\alpha$ and slope $\beta$) are estimated from the sample data. For $n = 28$:
  \[ df = 28 - 2 = 26 \]

  \item \textbf{Answer: (A)} \\
  \textbf{Explanation:} An appropriate linear model displays a random cloud of points scattered evenly above and below the zero residual line with constant vertical spread. Any clear curved or parabolic pattern (such as a U-shape) in the residual plot indicates that the relationship between $x$ and $y$ is non-linear and that a straight line is not an appropriate model.

  \item \textbf{Answer: (B)} \\
  \textbf{Explanation:} Let $X$ be the number of rolls until a 4 appears ($p = 1/6$). The event that $X > 3$ means that the first three consecutive rolls must all be failures (not a 4). The probability of failure on any single roll is $q = 1 - 1/6 = 5/6$. By the multiplication rule for independent trials:
  \[ P(X > 3) = (5/6)^3 = \frac{125}{216} \approx 0.5787 \]

  \item \textbf{Answer: (B)} \\
  \textbf{Explanation:} When sampling without replacement from a finite population, the observations are technically dependent. However, if the sample size $n$ does not exceed 10\% of the population size $N$ ($n \le 0.10N$), the dependence between trials is negligible and the standard deviation formula $\sigma_{\hat{p}} = \sqrt{\frac{p(1-p)}{n}}$ remains sufficiently accurate.

  \item \textbf{Answer: (A)} \\
  \textbf{Explanation:} Since the entire 99\% confidence interval $(1.2, 5.8)$ is strictly greater than zero, zero is not a plausible value for the true difference $\mu_1 - \mu_2$. This provides statistically significant evidence at $\alpha = 0.01$ that $\mu_1 - \mu_2 > 0$, or $\mu_1 > \mu_2$.

  \item \textbf{Answer: (C)} \\
  \textbf{Explanation:} Power is the probability of correctly rejecting a false null hypothesis ($1 - \beta$). Increasing the sample size $n$ decreases the standard error of the estimator, narrowing the sampling distribution and reducing the overlap between the null and alternative distributions, which directly increases power.

  \item \textbf{Answer: (B)} \\
  \textbf{Explanation:} In a balanced completely randomized design with $N = 60$ subjects divided equally among $k = 3$ treatment groups, each group receives:
  \[ \frac{N}{k} = \frac{60}{3} = 20\text{ subjects} \]

  \item \textbf{Answer: (B)} \\
  \textbf{Explanation:} In a normal distribution, the Empirical Rule states that approximately 68\% of values lie within $\pm 1\sigma$ of the mean, leaving 32\% in the two tails (16\% below $\mu - 1\sigma$ and 16\% above $\mu + 1\sigma$). Therefore, the area below $\mu + 1\sigma$ is $50\% + 34\% = 84\%$, corresponding to the $84^{\text{th}}$ percentile:
  \[ x = \mu + 1\sigma = 500 + 100(1) = 600 \]

  \item \textbf{Answer: (C)} \\
  \textbf{Explanation:} Standard deviation $s_x$ is computed in the exact same measurement units as the original data (e.g., inches, grams, dollars), making it readily interpretable. Variance is measured in squared units. Standard deviation is non-resistant to outliers, can never be negative, and scales linearly when multiplied by a constant.

  \item \textbf{Answer: (B)} \\
  \textbf{Explanation:} A fundamental mathematical property of the least-squares regression line (LSRL) $\hat{y} = b_0 + b_1 x$ is that it always passes through the point of averages, $(\bar{x}, \bar{y})$. This is derived from the formula for the intercept: $b_0 = \bar{y} - b_1 \bar{x}$.

  \item \textbf{Answer: (B)} \\
  \textbf{Explanation:} Voluntary response bias occurs when sample members self-select themselves into the survey by choosing whether to respond (such as callers responding to a call-in radio poll or online survey). People with strong negative opinions are heavily overrepresented.

  \item \textbf{Answer: (C)} \\
  \textbf{Explanation:} For two independent events $A$ and $B$, the multiplication rule states:
  \[ P(A \cap B) = P(A) \times P(B) = 0.3 \times 0.7 = 0.21 \]

  \item \textbf{Answer: (B)} \\
  \textbf{Explanation:} The $p$-value is the conditional probability, assuming the null hypothesis $H_0$ is true, of observing a test statistic as extreme as or more extreme than the value actually observed in the sample, in the direction of the alternative hypothesis $H_a$.

  \item \textbf{Answer: (D)} \\
  \textbf{Explanation:} For a two-way contingency table with $r$ rows and $c$ columns, the degrees of freedom for a Chi-Square test of independence or homogeneity are:
  \[ df = (r - 1)(c - 1) = (2 - 1)(3 - 1) = 1 \times 2 = 2 \]

  \item \textbf{Answer: (A)} \\
  \textbf{Explanation:} The slope of the least-squares regression line is calculated as:
  \[ b_1 = r \left(\frac{s_y}{s_x}\right) = -0.6 \left(\frac{10}{5}\right) = -0.6(2) = -1.2 \]

  \item \textbf{Answer: (B)} \\
  \textbf{Explanation:} The Large Counts condition for Chi-Square tests requires that all individual expected cell counts must be at least 5 ($E_i \ge 5$). This ensures that the discrete multinomial distribution is sufficiently well approximated by the continuous Chi-Square distribution.

  \item \textbf{Answer: (B)} \\
  \textbf{Explanation:} The 1-sample $t$-test statistic is:
  \[ t = \frac{\bar{x} - \mu_0}{s / \sqrt{n}} = \frac{50 - 45}{10 / \sqrt{25}} = \frac{5}{10 / 5} = \frac{5}{2} = 2.5 \]

  \item \textbf{Answer: (A)} \\
  \textbf{Explanation:} In a double-blind experiment, neither the experimental subjects nor the researchers/evaluators who interact with them know which treatment each subject is receiving. This eliminates both the psychological placebo effect in subjects and diagnostic/evaluator bias in researchers.

  \item \textbf{Answer: (C)} \\
  \textbf{Explanation:} The sample proportion $\hat{p}$ is an unbiased estimator of the true population proportion $p$. The mean of its sampling distribution is exactly equal to the parameter: $\mu_{\hat{p}} = p$.

  \item \textbf{Answer: (B)} \\
  \textbf{Explanation:} For the standard normal distribution $Z \sim N(0, 1)$, the area to the left of $z = -1.96$ is:
  \[ P(Z \le -1.96) = 0.0250 \]
  Together with the right tail above $+1.96$ ($0.0250$), exactly 5\% of the total area lies in the two tails, leaving 95\% in the center.

  \item \textbf{Answer: (A)} \\
  \textbf{Explanation:} The five conditions required for linear regression inference are summarized by the acronym \textbf{LINER}: Linear relationship, Independent observations, Normal distribution of $y$-values for each $x$, Equal variance of residuals ($\sigma$ constant), and Random sampling or assignment.

  \item \textbf{Answer: (B)} \\
  \textbf{Explanation:} Because each subject receives both conditions (resting and post-exercise measurements) and the analysis focuses on the within-subject differences, this is an example of a matched pairs experimental design.

  \item \textbf{Answer: (B)} \\
  \textbf{Explanation:} By the complement rule of probability:
  \[ P(A^c) = 1 - P(A) = 1 - 0.8 = 0.2 \]

  \item \textbf{Answer: (C)} \\
  \textbf{Explanation:} By mathematical construction of the ordinary least-squares line, the positive and negative vertical deviations from the line perfectly cancel out: $\sum (y_i - \hat{y}_i) = 0$.

  \item \textbf{Answer: (B)} \\
  \textbf{Explanation:} The width of the confidence interval is $\text{Upper} - \text{Lower} = 18.6 - 12.4 = 6.2$. The margin of error is half the interval width:
  \[ \text{ME} = \frac{6.2}{2} = 3.1 \]

  \item \textbf{Answer: (C)} \\
  \textbf{Explanation:} The number of cars passing an intersection is obtained by counting distinct individual vehicles ($0, 1, 2, 3, \dots$). Values are non-negative integers with no fractional values possible, defining a discrete quantitative variable.

  \item \textbf{Answer: (A)} \\
  \textbf{Explanation:} In hypothesis testing, the decision rule is to reject the null hypothesis $H_0$ if and only if the $p\text{-value} \le \alpha$. Therefore, if $H_0$ was rejected at $\alpha = 0.05$, the $p$-value must be less than or equal to 0.05.

  \item \textbf{Answer: (B)} \\
  \textbf{Explanation:} Margin of error is $\text{ME} = z^* \left(\frac{\sigma}{\sqrt{n}}\right)$. Increasing the confidence level from 90\% ($z^* = 1.645$) to 99\% ($z^* = 2.576$) increases the critical value $z^*$, which directly increases the margin of error and widens the confidence interval.

  \item \textbf{Answer: (B)} \\
  \textbf{Explanation:} The standard error of the slope, $\text{SE}(b_1)$, estimates the standard deviation of the sampling distribution of sample slopes $b_1$ across repeated random samples of size $n$ drawn from the same population.

  \item \textbf{Answer: (B)} \\
  \textbf{Explanation:} The coefficient of determination is $r^2$. Given the correlation coefficient $r = -0.68$:
  \[ r^2 = (-0.68)^2 = 0.4624 \implies 46.2\% \]
  Interpretation: Approximately 46.2\% of the variability in resting pulse rate is accounted for by the linear relationship with daily physical activity. Option (C) 68.0\% confuses $r$ with $r^2$.

  \item \textbf{Answer: (A)} \\
  \textbf{Explanation:} In simple linear regression, the test statistic for the slope is:
  \[ t = \frac{b_1 - 0}{\text{SE}(b_1)} = \frac{-0.32}{0.075} \approx -4.267 \approx -4.27 \]
  Degrees of freedom for inference on slope are $df = n - 2 = 25 - 2 = 23$. Option (B) incorrectly uses $df = n - 1 = 24$.
\end{enumerate}

\section{Section II: Free-Response Complete Scoring Guidelines}

\begin{frqrubricbox}[Question 1: Multi-Focus on Practices 1 \& 2 --- Model Solution \& Rubric]
\textbf{Part (a) Skill 1.A: Formulate Question}\\
\textit{"Does replacing traditional seated desks with dynamic standing desks lead to a statistically significant increase in the mean academic attentiveness rating of high school students?"}

\textbf{Part (b) Skill 2.B: Random Assignment Protocol}\\
1. Number the 40 available classrooms sequentially from 01 to 40.\\
2. Using a random number generator or statistical software, generate 20 unique pseudo-random integers between 01 and 40 without replacement.\\
3. The classrooms corresponding to these 20 numbers will receive the dynamic standing desks (Treatment Group). The remaining 20 classrooms will retain traditional seated desks (Control Group).\\
4. Following an 8-week intervention period, record attentiveness using an identical standardized rubric.\\
5. Compare the mean attentiveness scores between the two treatment conditions.

\textbf{Part (c) Skill 2.D: Confounding \& Randomized Block Design}\\
Academic subject matter (e.g., Mathematics vs. Visual Arts) is a potential confounding variable because student attentiveness and natural physical movement vary systematically across subjects. If one desk type happens to be assigned to more physical education or art classes, the effect of desk type is confounded with academic subject matter.\\
\textit{Blocking Remedy:} Group the 40 classrooms into homogeneous blocks by academic discipline (e.g., Block 1: STEM classes; Block 2: Humanities/Arts). Within each block, randomly assign half the classrooms to standing desks and half to traditional desks. This eliminates between-subject variability from the experimental error.

\textbf{Part (d) Skill 2.E: Null \& Alternative Hypotheses}\\
Let $\mu_S$ be the true mean attentiveness score of classrooms with standing desks, and let $\mu_T$ be the true mean attentiveness score of classrooms with traditional desks:
\[ H_0: \mu_S = \mu_T \quad (\text{or } \mu_S - \mu_T = 0) \]
\[ H_a: \mu_S > \mu_T \quad (\text{or } \mu_S - \mu_T > 0) \]

\textbf{Grading Rubric Breakdown:}\\
\begin{itemize}
  \item \textbf{Essentially Correct (E):} Clearly formulates an investigative comparative question with variable and population, provides a fully reproducible random allocation procedure, thoroughly explains confounding with a valid blocking design, and states hypotheses correctly defining parameters in context.
  \item \textbf{Partially Correct (P):} Omits sampling without replacement in part (b), or identifies confounding without explaining how blocking isolates the variation, or writes hypotheses using sample statistics ($\bar{x}_S = \bar{x}_T$).
  \item \textbf{Incomplete (I):} Major conceptual confusion between blocking and stratified sampling.
\end{itemize}
\end{frqrubricbox}

\begin{frqrubricbox}[Question 2: Multi-Focus on Practices 3 \& 4 --- Model Solution \& Rubric]
\textbf{Part (a) Skill 3.B: Outlier Identification}\\
Ordered sugar values ($n = 14$): $4, 6, 7, 8, 9, 10, 11, 12, 12, 13, 14, 15, 18, 26$.\\
$\text{Median} = (11 + 12)/2 = 11.5\text{ g}$. Lower half (first 7): $Q_1 = 8\text{ g}$. Upper half (last 7): $Q_3 = 14\text{ g}$.\\
$\text{IQR} = Q_3 - Q_1 = 14 - 8 = 6\text{ g}$.\\
$\text{Lower Fence} = 8 - 1.5(6) = 8 - 9 = -1\text{ g}$.\\
$\text{Upper Fence} = 14 + 1.5(6) = 14 + 9 = 23\text{ g}$.\\
Because $26\text{ g} > 23\text{ g}$, 26 grams is a confirmed high outlier. There are no low outliers ($\text{Min} = 4 \ge -1$).

\textbf{Part (b) Skill 4.A: C-U-S-S Distribution Description}\\
\textit{"The distribution of sugar content among these 14 breakfast cereals is skewed to the right with an isolated high outlier at 26 grams per serving. The center of the distribution is represented by a median of 11.5 grams (or mean of approximately 12.07 grams). The spread is characterized by an IQR of 6 grams and an overall range of 22 grams (from 4 g to 26 g)."}

\textbf{Part (c) Skill 3.B: Prediction \& Residual Calculation}\\
For $x = 12.0\text{ g}$: $\widehat{\text{Calories}} = 95.20 + 3.85(12.0) = 95.20 + 46.20 = 141.40\text{ calories}$.\\
$\text{Residual} = y - \hat{y} = 148.0 - 141.40 = +6.60\text{ calories}$.\\
\textit{Interpretation:} The cereal contains 6.60 more calories per serving than predicted by the linear regression model based on its sugar content.

\textbf{Part (d) Skill 4.B: Parameter \& Model Interpretations}\\
\textit{Slope (3.85):} For each additional 1 gram increase in sugar content per serving, the model predicts that caloric content increases by approximately 3.85 calories.\\
\textit{Coefficient of Determination ($r^2 = 72.8\%$):} Approximately 72.8\% of the total variation in caloric content among these cereals is accounted for by the linear model with sugar content.

\textbf{Grading Rubric Breakdown:}\\
\begin{itemize}
  \item \textbf{Essentially Correct (E):} Correct 5-number summary and mathematical justification of outlier, complete C-U-S-S description in context, correct prediction with positive residual interpretation, and flawless slope/$r^2$ interpretations with context and units.
  \item \textbf{Partially Correct (P):} Forgets units in residual or slope interpretation, or omits one C-U-S-S component.
  \item \textbf{Incomplete (I):} Arithmetic errors leading to incorrect fences or states slope deterministically ("calories will always increase by 3.85").
\end{itemize}
\end{frqrubricbox}

\begin{frqrubricbox}[Question 3: Statistical Inference (Hypothesis Test) --- Model Solution \& Rubric]
\textbf{Step 1 (P): Parameter \& Hypotheses}\\
Let $p$ be the true proportion of all admitted high school seniors who accept their admission offer and enroll at the university.
\[ H_0: p = 0.65 \quad \text{vs.} \quad H_a: p < 0.65 \]

\textbf{Step 2 (H \& A): Name Procedure \& Check Conditions}\\
Procedure: One-proportion $z$-test for $p$.
\begin{itemize}
  \item \textbf{Random:} An SRS of $n = 300$ admitted students was selected. (Satisfied)
  \item \textbf{10\% Condition:} $300 \le 0.10(\text{Total cohort of admitted students})$, assuming the university admitted at least 3,000 students. (Satisfied)
  \item \textbf{Large Counts:} Under $H_0$, $n p_0 = 300(0.65) = 195 \ge 10$ and $n(1 - p_0) = 300(0.35) = 105 \ge 10$. Both expected counts exceed 10. (Satisfied)
\end{itemize}

\textbf{Step 3 (N \& T): Mechanics \& Calculations}\\
Sample proportion: $\hat{p} = \frac{180}{300} = 0.60$.
Standard error under $H_0$:
\[ \text{SE}_0 = \sqrt{\frac{p_0(1 - p_0)}{n}} = \sqrt{\frac{(0.65)(0.35)}{300}} = \sqrt{\frac{0.2275}{300}} \approx 0.02754 \]
Standardized test statistic:
\[ z = \frac{\hat{p} - p_0}{\text{SE}_0} = \frac{0.60 - 0.65}{0.02754} = \frac{-0.05}{0.02754} \approx -1.82 \]
$p$-value: $P(Z \le -1.82) = 0.0344$.

\textbf{Step 4 (O \& M): Decision \& Contextual Conclusion}\\
Because the $p$-value ($0.0344$) is strictly less than the significance level $\alpha = 0.05$, we \textbf{reject the null hypothesis $H_0$}. There is convincing statistical evidence at the 5\% level that the true enrollment proportion among admitted students is less than 65\%.

\textbf{Grading Rubric Breakdown:}\\
\begin{itemize}
  \item \textbf{Essentially Correct (E):} Correct parameter definition, all 3 conditions verified with numbers, correct test statistic $z = -1.82$ and $p$-value $0.0344$, and correct contextual conclusion explicitly comparing $p$-value to $\alpha$.
  \item \textbf{Partially Correct (P):} Uses $\hat{p}$ instead of $p_0$ in standard error formula or states hypotheses without defining $p$.
  \item \textbf{Incomplete (I):} Omits conditions or makes an incorrect decision ("accept $H_0$").
\end{itemize}
\end{frqrubricbox}

\begin{frqrubricbox}[Question 4: Multi-Focus Synthesis on Practices 2, 3, \& 4 --- Model Solution \& Rubric]
\textbf{Part (a): Residual \& Slope Interpretation}\\
Predicted score for $x = 4.0$: $\hat{y} = 42.0 + 7.5(4.0) = 42.0 + 30.0 = 72.0\text{ points}$.\\
$\text{Residual} = y - \hat{y} = 78.0 - 72.0 = +6.0\text{ points}$.\\
\textit{Interpretation:} The student scored 6.0 points higher on the exam than predicted by study hours alone. The slope indicates that for each additional hour of daily study, the predicted exam score increases by 7.5 points.

\textbf{Part (b): Omitted Variable Bias Analysis}\\
In the multiple regression model, the slope for study hours drops from 7.5 to 6.0. The simple linear regression model suffered from \textbf{omitted variable bias}: students who study more may also manage their sleep differently, or sleep duration ($w$) independently explains score variability. Omitting sleep attributed some of the collective gains to study time alone, overestimating its marginal effect.

\textbf{Part (c): Probability Model Calculation}\\
$G \sim N(12.5, 3.2)$. Standardizing $g = 15.0$:
\[ z = \frac{15.0 - 12.5}{3.2} = \frac{2.5}{3.2} \approx 0.78 \]
\[ P(G > 15.0) = P(Z > 0.78) = 1 - 0.7823 = 0.2177 \quad (\text{or } 21.8\%) \]
Interpretation: Approximately 21.8\% of students are expected to achieve score gains exceeding 15.0 points.

\textbf{Part (d): Constrained Optimization \& Policy Recommendation}\\
Given constraint $x + w \le 14 \implies w = 14 - x$, with $w \ge 6 \implies x \le 8$.
Substitute into multiple regression equation:
\[ \hat{y} = 20.0 + 6.0x + 3.5(14 - x) = 20.0 + 6.0x + 49.0 - 3.5x = 69.0 + 2.5x \]
Because the effective derivative with respect to study time is $+2.5 > 0$, predicted score is maximized by choosing the largest feasible value of $x$. Subject to minimum sleep $w \ge 6$, the maximum allowable study time is $x = 14 - 6 = 8.0\text{ hours}$.
Maximum predicted score:
\[ \hat{y}_{\max} = 69.0 + 2.5(8) = 69.0 + 20.0 = 89.0\text{ points} \]
\textit{Recommendation:} Students should allocate 8 hours to studying and 6 hours to sleep to maximize predicted test scores under the 14-hour daily budget.

\textbf{Grading Rubric Breakdown:}\\
\begin{itemize}
  \item \textbf{Essentially Correct (E):} Correct residual calculation with context, sound conceptual explanation of omitted variable bias, accurate normal distribution probability with work shown, and mathematically rigorous linear optimization under constraints.
  \item \textbf{Partially Correct (P):} Correctly calculates parts (a) and (c) but fails to incorporate the boundary constraint $w \ge 6$ in part (d).
  \item \textbf{Incomplete (I):} Major misunderstandings in regression mechanics or normal probability.
\end{itemize}
\end{frqrubricbox}
